\chapter{李群、李代数与齐性空间}
\label{ch:15-lie-groups}
\ChapterPlan{周第5章（本地页码待核对）；Kirillov讲义第2--4、6章作补充}{正文证明初稿；8题详解与3幅原创图，尚未独立数学审阅}

把一个几何变换连续地改变，便同时遇到群的乘法与流形的微分。本章从单位元处的切空间恢复局部乘法的无穷小信息，再说明闭子群如何产生光滑商空间。最后把不变度量与测地线反射联系起来。所有李群均指有限维、Hausdorff、第二可数、无边界的光滑实李群；不预设连通。李代数括号沿用第5章的导子交换子，曲率沿用第11章的符号。

证明依赖为：第5章的ODE及常秩定理用于指数映射和作用，第8章的积分用于紧群平均，第11--12章的Levi--Civita联络与Jacobi场用于不变度量和局部对称性。闭子群定理及闭子群的商流形在本章补证，不作为隐藏先修。Lie第三定理、任意子代数的浸入积分、完整BCH公式、根系分类和全体等距群的Lie群结构不在本章证明范围内。

\section{李群、左不变向量场与李代数}
\subsection{群运算的微分与基本例子}
\begin{definition}
李群是一个群$G$及其光滑流形结构，使乘法$\mu(g,h)=gh$与取逆$\iota(g)=g^{-1}$光滑。单位元记$e$，左、右平移分别为$L_g(h)=gh$、$R_g(h)=hg$。李群同态是保持乘法的光滑映射。
\end{definition}
每个$L_g,R_g$均为微分同胚，逆分别为$L_{g^{-1}},R_{g^{-1}}$。令$\mathfrak g=T_eG$。由$\mu(g,e)=g$及$\mu(e,h)=h$，微分的线性性给出
\[
 \dd\mu_{(e,e)}(X,Y)=X+Y.
\]
再微分$g\iota(g)=e$，得到$\dd\iota_e=-I$。这两条公式只在单位元作这种简单表示；不能把一般群乘法看成切空间中的加法。

\begin{example}[加法群、圆周与矩阵群]
$\R^n$的加法与$S^1\subset\mathbb C$的复数乘法均为李群。$\mathrm{GL}(n,\R)$是矩阵空间中$\det\ne0$的开集；逆矩阵由余子式除以行列式表示，故取逆光滑。

在可逆矩阵$A$处，行列式的一阶项为
\[
 \dd(\det)_A(V)=\det(A)\operatorname{tr}(A^{-1}V).
\]
这是由$\det(A+tV)=\det A\det(I+tA^{-1}V)$及行列式展开中一次项等于迹得到的。该导数满射到$\R$，所以$\mathrm{SL}(n,\R)=\{\det A=1\}$是正则水平集，其单位元切空间为迹零矩阵。

对$F(A)=A^{\mathsf T}A$，在$A\in\mathrm O(n)$处有
\[
 \dd F_A(V)=V^{\mathsf T}A+A^{\mathsf T}V.
\]
余域取对称矩阵空间；给定对称$S$，选$V=AS/2$即映为$S$。因此$\mathrm O(n)=F^{-1}(I)$为嵌入子流形，切空间由$A^{\mathsf T}V$反对称刻画。$\mathrm{SO}(n)$是在其中$\det=1$的开闭子流形。群运算是环境运算的限制，于是这些都是李群，且
\[
 \mathfrak{gl}(n,\R)=M_n(\R),\quad
 \mathfrak{sl}(n,\R)=\{X:\operatorname{tr}X=0\},\quad
 \mathfrak{so}(n)=\{X:X^{\mathsf T}=-X\}.
\]
此处已给出具体子流形证明，不依赖稍后的一般闭子群定理。
\end{example}

\subsection{左不变场与括号}
\begin{proposition}\label{prop:lie-left-fields}
对$X\in\mathfrak g$，公式$X^L(g)=\dd(L_g)_eX$定义光滑左不变场。它给出$\mathfrak g$与全部左不变场的线性同构。定义
\[
 [X,Y]=[X^L,Y^L](e),
\]
则$[X^L,Y^L]=[X,Y]^L$，且此括号满足反对称性与Jacobi恒等式，称为$G$的李代数。
\end{proposition}
\begin{proof}
光滑性由群乘法的微分给出。$L_hL_g=L_{hg}$说明左不变性。反过来，左不变场$V$必满足$V(g)=\dd(L_g)_eV(e)$，故单位元取值确定全场。微分同胚保持向量场的括号，而$L_h$保持$X^L,Y^L$，所以也保持它们的括号。反对称性及Jacobi恒等式因而从第5章的向量场括号传到单位元。
\end{proof}

对矩阵群，$X^L(A)=AX$。环境线性场的导数给出
\[
 [X^L,Y^L](A)=(AX)Y-(AY)X=A(XY-YX),
\]
故矩阵李代数的括号确实是$[X,Y]=XY-YX$。迹的循环性表明迹零矩阵封闭；反对称矩阵的交换子仍反对称，也与上面的切空间相容。

右不变场记$X^R(g)=\dd(R_g)_eX$。取逆映射满足
$\iota\circ L_g=R_{g^{-1}}\circ\iota$，从而$\iota_*X^L=-X^R$。将括号也推过去，得到必须记住的反号
\begin{equation}\label{eq:lie-right-sign}
 [X^R,Y^R]=-[X,Y]^R.
\end{equation}
例如矩阵群的右不变场为$X^R(A)=XA$，直接求导便是$(YX-XY)A$。后面的群作用也将遇到这个反号。


\Needspace{17\baselineskip}
\section{指数映射与群同态}
\subsection{一参数子群与完备性}
\begin{theorem}\label{thm:lie-exponential}
每个左不变场$X^L$完备。存在唯一光滑曲线$\gamma_X:\R\to G$使$\gamma_X(0)=e$、$\dot\gamma_X=X^L\circ\gamma_X$，且
\[
 \gamma_X(s+t)=\gamma_X(s)\gamma_X(t),\qquad
 \gamma_{aX}(t)=\gamma_X(at).
\]
定义$\exp X=\gamma_X(1)$，则$\exp:\mathfrak g\to G$光滑，$\dd\exp_0=I$，故在零点附近是微分同胚。$X^L$的流为
\begin{equation}\label{eq:lie-left-flow}
 \Phi_t(g)=g\exp(tX)=R_{\exp(tX)}(g).
\end{equation}
\end{theorem}
\begin{proof}
先取单位元出发的局部解$\gamma:(-\varepsilon,\varepsilon)\to G$。对固定$s$，$t\mapsto\gamma(s)\gamma(t)$由左不变性也是$X^L$的积分曲线，初值为$\gamma(s)$；唯一性给出局部乘法式。把这条短解左移至任意已到达点，可在每个端点再延长同一个时间长度。反复拼接得到全实轴解，唯一性使重叠处一致，也把乘法式推广到一切$s,t$。时间缩放验证第二式。初值$g$的解便是$L_g\gamma_X$，得到\eqref{eq:lie-left-flow}。

参数ODE给出$X=0$附近一个固定短时间$\varepsilon>0$上的光滑解$\gamma_X(t)$。任取$X_0$，选足够大的整数$N$，使$X/(N\varepsilon)$对$X$在$X_0$的小邻域内落在这个局部参数范围。缩放及乘法式给出
\[
 \exp X=\bigl(\gamma_{X/(N\varepsilon)}(\varepsilon)\bigr)^N,
\]
从而得到该邻域的光滑性；右边只调用已经有参数光滑性的短时间解。最后$\exp(tX)=\gamma_X(t)$的零时导数为$X$，故$\dd\exp_0X=X$，逆函数定理给出结论。
\end{proof}

\begin{figure}[htbp]
\centering\DGDiagram{lie-left-flow}

\caption{左不变场的解从$e$经左平移搬到$g$；固定时间的流映射则是右平移。箭头注明了各自固定的量。}
\label{fig:lie-left-flow}
\end{figure}

任一光滑一参数子群$c:\R\to G$满足$c(s+t)=c(s)c(t)$，对$t$在零点求导得$c'(s)=\dd L_{c(s)}c'(0)$，所以必为$\exp(tX)$。右不变场的流相应为$g\mapsto\exp(tX)g$。

\begin{example}[矩阵指数]
矩阵级数$E(t)=\sum_{j\ge0}t^jX^j/j!$及逐项导数在有界$t$区间一致收敛，故$E'=EX$、$E(0)=I$。由ODE唯一性，它等于上面在任意矩阵李群中定义的$\exp(tX)$；特别地解始终留在该群。$S^1$的指数为$t\mapsto e^{it}$，显然不单射。局部可逆不能推出全局单射或满射，习题15.6将给出连通单位分支中的非指数元素。
\end{example}

\subsection{同态的微分、刚性与覆盖}
\begin{theorem}\label{thm:lie-homomorphism}
若$F:G\to K$为光滑群同态，$f=\dd F_e$，则$f$是李代数同态，并且
\[
 F(\exp_G X)=\exp_K(fX),\qquad
 \dd F_g\circ\dd L_g=\dd L_{F(g)}\circ f.
\]
$F$常秩，$\ker F$是闭嵌入正规李子群，$T_e\ker F=\ker f$。若$G$连通，则$F$由$f$唯一确定。
\end{theorem}
\begin{proof}
第二式来自微分$F\circ L_g=L_{F(g)}\circ F$，其左右平移微分可逆，故秩常数。它也说明$X^L$与$(fX)^L$是$F$相关的场；相关场的括号相关，单位元取值给出$f[X,Y]=[fX,fY]$。将积分曲线用$F$映出，由ODE唯一性得到指数的自然性。常秩定理把单位纤维局部写成坐标子空间，故核是嵌入子流形，切空间是微分的核。核闭且正规来自连续性与群同态恒等式。

连通群由任意单位元邻域$U$生成：$U\cup U^{-1}$的有限乘积之并是开子群，所有陪集也开，故其补集开；连通性迫使它为全群。两个同态若微分相等，便由指数自然性在一个单位元邻域相等，继而在其生成的全群相等。连通性不可删去，例如离散群的所有切空间都为零。
\end{proof}

\Needspace{12\baselineskip}
\begin{proposition}\label{prop:lie-covering}
若$f$为线性同构且$K$连通，则$F$满射，是具有离散核的覆盖映射。若$G$也连通，则核包含在$G$的中心。
\end{proposition}
\begin{proof}
逆函数定理给出单位元邻域$U$，使$F|_U:U\to V$微分同胚。像包含$V$，故是开子群；$K$连通使像为全体。缩小$U$使$UU^{-1}\cap\ker F=\{e\}$。若$F(g)\in V$，存在唯一$u\in U$满足$F(u)=F(g)$，于是$g=ku$且$k\in\ker F$。不同$kU$互不相交，因为$ku=\ell v$蕴含$\ell^{-1}k=vu^{-1}\in UU^{-1}\cap\ker F$。所以$F^{-1}(V)$是这些开集的不交并，每片映到$V$均为微分同胚；平移给出所有点的均匀覆盖邻域。最后对$k\in\ker F$，映射$g\mapsto gkg^{-1}$从连通$G$进入离散核，必恒为$k$。
\end{proof}

这里没有声称每个李代数同态都能在任意给定的源群上积分。例如$\R\to\R$的任意线性映射是李代数同态，但要下降为$S^1\to S^1$同态，必须保持指数的周期格；习题15.4会算出整数条件。


\Needspace{17\baselineskip}
\section{伴随表示与Killing形式}
\subsection{伴随表示与Maurer--Cartan方程}
\begin{definition}
共轭$c_g(h)=ghg^{-1}$的单位元微分称为伴随作用$\operatorname{Ad}_g:\mathfrak g\to\mathfrak g$。由于$c_g$为群自同构，它是李代数自同构，且$\operatorname{Ad}_{gh}=\operatorname{Ad}_g\operatorname{Ad}_h$。
\end{definition}
\begin{proposition}\label{prop:lie-ad}
令$\operatorname{ad}_X(Y)=[X,Y]$，则
\[
 \left.\frac{\dd}{\dd t}\right|_0\operatorname{Ad}_{\exp(tX)}Y=[X,Y],\qquad
 \operatorname{Ad}_{\exp(tX)}=e^{t\operatorname{ad}_X}.
\]
\end{proposition}
\begin{proof}
第5章的流公式为$[V,W]=\left.\frac{\dd}{\dd t}\right|_0(\Phi_{-t})_*W$。对$V=X^L,W=Y^L$，在$e$处取值，$\Phi_t=R_{\exp(tX)}$，故被求导的向量为
\[
 \dd(R_{\exp(-tX)})_{\exp(tX)}\dd(L_{\exp(tX)})_eY
 =\operatorname{Ad}_{\exp(tX)}Y.
\]
这给出第一式。令$A(t)=\operatorname{Ad}_{\exp(tX)}$，一参数乘法式使$A'=A\operatorname{ad}_X$、$A(0)=I$，解为第二式。Jacobi恒等式也直接给出$[\operatorname{ad}_X,\operatorname{ad}_Y]=\operatorname{ad}_{[X,Y]}$。
\end{proof}

\begin{proposition}[左Maurer--Cartan形式]\label{prop:lie-maurer-cartan}
定义$\mathfrak g$值一形式$\theta_g=\dd(L_{g^{-1}})_g$，则$\theta(X^L)=X$且
\begin{equation}\label{eq:lie-maurer-cartan}
 \dd\theta+\tfrac12[\theta\wedge\theta]=0,
 \qquad [\theta\wedge\theta](u,v)=2[\theta(u),\theta(v)].
\end{equation}
其中外微分按$\mathfrak g$任一固定基逐分量定义。矩阵群中$\theta=g^{-1}\dd g$。
\end{proposition}
\begin{proof}
对左不变场，$\dd\theta(X^L,Y^L)=X^L(Y)-Y^L(X)-\theta([X^L,Y^L])=-[X,Y]$。它恰好与第二项抵消；左不变场逐点张成切空间，二形式的张量性遂给出全式。矩阵表达来自左平移的微分。
\end{proof}

\subsection{Killing形式及其限度}
\begin{definition}
李代数上的Killing形式为$B(X,Y)=\operatorname{tr}(\operatorname{ad}_X\operatorname{ad}_Y)$。它对称，但未必正定、负定或非退化。
\end{definition}
\begin{proposition}\label{prop:lie-killing}
$B$在每个李代数自同构下不变，且
\[
 B([Z,X],Y)+B(X,[Z,Y])=0.
\]
其根空间$\{X:B(X,Y)=0\text{ 对所有 }Y\}$为理想。
\end{proposition}
\begin{proof}
自同构$A$满足$\operatorname{ad}_{AX}=A\operatorname{ad}_XA^{-1}$，迹在共轭下不变给出第一条。利用上一命题的交换子式，第二条的左边是
\[
 \operatorname{tr}\bigl([\operatorname{ad}_Z,\operatorname{ad}_X]\operatorname{ad}_Y
 +\operatorname{ad}_X[\operatorname{ad}_Z,\operatorname{ad}_Y]\bigr)
 =\operatorname{tr}[\operatorname{ad}_Z,\operatorname{ad}_X\operatorname{ad}_Y]=0.
\]
若$X$在根空间中，则$B([Z,X],Y)=-B(X,[Z,Y])=0$，故该空间为理想。
\end{proof}

\begin{example}
交换李代数的所有$\operatorname{ad}$均为零，故$B=0$。把$\R^3$通过$\widehat v(w)=v\times w$识别为$\mathfrak{so}(3)$，向量三重积给出$[\widehat v,\widehat w]=\widehat{v\times w}$，且
\[
 B(\widehat v,\widehat w)=\operatorname{tr}(\widehat v\widehat w)=-2\langle v,w\rangle.
\]
对$\mathfrak{sl}(2,\R)$取
$H=\left(\begin{smallmatrix}1&0\\0&-1\end{smallmatrix}\right)$、
$E=\left(\begin{smallmatrix}0&1\\0&0\end{smallmatrix}\right)$、
$F=\left(\begin{smallmatrix}0&0\\1&0\end{smallmatrix}\right)$，有
$[H,E]=2E,[H,F]=-2F,[E,F]=H$。逐基计算得$B(H,H)=8$、$B(E,F)=4$，其余基配对为零，即$B(X,Y)=4\operatorname{tr}(XY)$。所以$H,E+F$为正方向，$E-F$为负方向；不能把$-B$总当黎曼内积。这里不以未证明的Cartan判据推断一般李代数的结构。
\end{example}

\subsection{不变度量与紧群平均}
任取$\mathfrak g$上的正定内积$\langle\ ,\ \rangle$，通过左平移唯一延拓为左不变黎曼度量。把$\dd R_h$用左平凡化写出，得到
\[
 \theta_{gh}(\dd R_h v)=\operatorname{Ad}_{h^{-1}}\theta_g(v).
\]
因此该度量右不变当且仅当内积在全部$\operatorname{Ad}(G)$下不变，称为双不变度量。若$G$连通，等价条件为
\begin{equation}\label{eq:lie-ad-skew}
 \langle[X,Y],Z\rangle+\langle Y,[X,Z]\rangle=0.
\end{equation}
必要性由沿$\exp(tX)$求导得到；充分性由$e^{t\operatorname{ad}_X}$正交、指数邻域生成连通群得到。非连通时还须检查其余分支的$\operatorname{Ad}$。

\begin{theorem}\label{thm:lie-compact-average}
紧李群具有双不变黎曼度量。紧子群作用在有限维实向量空间上时，也存在使作用正交的正定内积。
\end{theorem}
\begin{proof}
先在$G$上选一个左不变正密度$\nu$：取$T_eG$上的正密度，以左平移延拓。第8章的积分在各坐标中对正密度同样成立，不需定向。因左右平移交换，$R_h^*\nu$仍左不变，故$R_h^*\nu=c(h)\nu$，其中$c:G\to\R_{>0}$为连续群同态（乘法次序在标量群中无影响）。紧子群$c(G)$只能是$\{1\}$：若含$a\ne1$，则$a$或$a^{-1}$大于1，其幂无界。因此$\nu$也右不变。紧致性使总质量有限且非零，归一化为1。

给定初始正定内积$b_0$，令
\[
 b(X,Y)=\int_G b_0(\operatorname{Ad}_gX,\operatorname{Ad}_gY)\,\nu(g).
\]
这是正定双线性型，右不变换元$g\mapsto gh$给出$\operatorname{Ad}_h$不变性。一般表示$\rho:H\to\mathrm{GL}(V)$以$\rho(h)$代替$\operatorname{Ad}_g$作同一平均即可。这里证明了紧李群所需的积分平均，没有调用一般局部紧群的Haar测度定理。
\end{proof}

\begin{theorem}\label{thm:lie-biinvariant-curvature}
对双不变度量及左不变场，有
\begin{align*}
 \nabla_{X^L}Y^L&=\tfrac12[X,Y]^L,\\
 R(X,Y)Z&=-\tfrac14[[X,Y],Z],\\
 \langle R(X,Y)Y,X\rangle&=\tfrac14\|[X,Y]\|^2,\qquad
 \operatorname{Ric}(X,Y)=-\tfrac14B(X,Y).
\end{align*}
故截面曲率非负，且$\exp(tX)$是从$e$出发的黎曼测地线。
\end{theorem}
\begin{proof}
左不变场两两内积为常数。Koszul公式及\eqref{eq:lie-ad-skew}给出$2\langle\nabla_XY,Z\rangle=\langle[X,Y],Z\rangle$。代入本书约定的$R=\nabla_X\nabla_Y-\nabla_Y\nabla_X-\nabla_{[X,Y]}$，Jacobi恒等式使前两项合为$\tfrac14[[X,Y],Z]$，减去最后的$\tfrac12[[X,Y],Z]$即得负号。再用内积不变性得到截面曲率的分子。对正交基$e_i$求迹，
\[
 \operatorname{Ric}(X,Y)=\tfrac14\sum_i\langle[e_i,X],[e_i,Y]\rangle
 =-\tfrac14\operatorname{tr}(\operatorname{ad}_X\operatorname{ad}_Y).
\]
最后沿$\exp(tX)$速度为$X^L$，协变加速度$\frac12[X,X]=0$。一般左不变度量不满足此结论，见习题15.8。
\end{proof}
此时$B(X,X)=-\|\operatorname{ad}_X\|_{\mathrm{HS}}^2$，所以Killing形式半负定，其根空间恰为中心。但反过来，仅凭半负定不能在未证明额外结构定理时便宣称群紧致。


\Needspace{17\baselineskip}
\section{齐性流形与黎曼对称空间}
\subsection{闭子群为何是子流形}
光滑同态的像未必是嵌入子流形；环面中的无理绕行将给出反例。以下定理的闭性条件正好排除局部累积出额外方向的现象。

\begin{theorem}[闭子群定理]\label{thm:lie-closed-subgroup}
李群$G$的任一闭子群$H$有唯一的嵌入子流形结构，使其为李子群。其李代数为
\begin{equation}\label{eq:lie-closed-algebra}
 \mathfrak h=\{X\in\mathfrak g:\exp(tX)\in H\text{ 对所有 }t\in\R\}.
\end{equation}
\end{theorem}
\begin{proof}
先定义$V$为所有经过$e$的$H$内$C^1$曲线在零点的速度。两条曲线相乘的速度是速度之和；重参数给出任意实数倍，所以$V$是有限维线性子空间。若$c'(0)=X$，在指数坐标中有$\log c(s)=sX+o(s)$。对固定$t$，当$n$充分大时
\[
 c(t/n)^n=\exp\bigl(n\log c(t/n)\bigr)\longrightarrow\exp(tX).
\]
左边属于$H$，闭性使极限属于$H$。反过来，$\exp(tX)$本身就是速度为$X$的曲线，故$V$恰为\eqref{eq:lie-closed-algebra}中的集合，记为$\mathfrak h$。它是线性子空间，因而闭。对$h\in H$，共轭曲线说明$\operatorname{Ad}_h\mathfrak h=\mathfrak h$；再对$h=\exp(tX)$求导，得$[X,Y]\in\mathfrak h$，所以它为李子代数。

关键的局部事实是：若$Z_j\ne0$、$Z_j\to0$、$\exp Z_j\in H$，且$Z_j/\|Z_j\|\to Z$，则$Z\in\mathfrak h$。事实上，对任意$t\in\R$选整数$n_j$使$n_j\|Z_j\|\to t$，有$(\exp Z_j)^{n_j}=\exp(n_jZ_j)\to\exp(tZ)$，再用闭性即可。

选补空间$\mathfrak g=\mathfrak m\oplus\mathfrak h$。映射
\begin{equation}\label{eq:lie-slice-product}
 \Psi(X,Y)=\exp X\exp Y,\qquad X\in\mathfrak m,\ Y\in\mathfrak h
\end{equation}
在$(0,0)$的微分为$(X,Y)\mapsto X+Y$，故在乘积小邻域上为微分同胚。还可缩小邻域，使$X\in\mathfrak m$且$\exp X\in H$蕴含$X=0$。否则能取这样的非零$X_j\to0$，由单位球紧性取归一化极限$Z\in\mathfrak m$，上一段又给出$Z\in\mathfrak h$，与$\|Z\|=1$矛盾。

于是$\Psi(X,Y)\in H$当且仅当$X=0$：其中$\exp Y\in H$，可用群乘法消去该因子。这在$e$附近把$H$精确写成坐标子空间，左平移到各点即为嵌入子流形。群运算的限制光滑，单位元切空间就是$\mathfrak h$。嵌入结构的光滑函数由环境局部坐标限制确定，故唯一。
\end{proof}

\subsection{左陪集的光滑商}
\begin{theorem}\label{thm:lie-quotient}
设$H$为闭子群，$G/H=\{gH:g\in G\}$取商拓扑；这里$gH$是左陪集，也是右乘$H$作用的轨道。它有唯一的光滑流形结构使$\pi:G\to G/H$为光滑满射下沉，并且
\[
 \dim(G/H)=\dim G-\dim H,\qquad
 T_{eH}(G/H)\simeq\mathfrak g/\mathfrak h.
\]
左乘给出$G$在$G/H$上的光滑传递作用。若$H$还正规，则商群为李群，其李代数为商李代数$\mathfrak g/\mathfrak h$。
\end{theorem}
\begin{proof}
$\pi$是开映射，因为$\pi^{-1}\pi(U)=\bigcup_{h\in H}Uh$开。若$gH\ne kH$，则$g^{-1}k\notin H$。闭性及乘法连续性给出$g,k$的小邻域$U,V$使$U^{-1}V$避开$H$；其像$\pi(U),\pi(V)$不交，故商Hausdorff。$G$的可数拓扑基经开映射$\pi$映出便是商的可数基。

在\eqref{eq:lie-slice-product}的乘积坐标中，$\pi\Psi(X,Y)=\pi(\exp X)$。先固定一个较大的可逆乘积坐标盒，再缩小$X$的邻域，使任意两个$X_1,X_2$对应的$\exp(-X_2)\exp X_1$都在$G$的上述坐标邻域内。若两者代表同一陪集，则
\[
 \exp(-X_2)\exp X_1\in H,
 \qquad \exp(-X_2)\exp X_1=\exp Y
\]
且$Y\in\mathfrak h$在原来的盒内。乘积坐标的唯一性用于$\exp X_1=\exp X_2\exp Y$，给出$X_1=X_2$。这也说明缩小邻域时必须同时控制横截面和子群方向。

故$X\mapsto\exp XH$是从$\mathfrak m$中开集到$\pi(\Psi(U\times V))$的连续双射。它也是开映射：对开$O\subset U$，$\Psi(O\times V)$在$G$中开，且其$\pi$像恰为$\{\exp XH:X\in O\}$。所以得到商的局部坐标。平移给出所有$gH$附近的坐标$X\mapsto g\exp XH$。

验证过渡光滑时，在重叠点固定$h_0\in H$，使$k^{-1}g\exp X_0h_0$位于$k$坐标所用的乘积邻域。对附近$X$，用$\Psi^{-1}$分解$k^{-1}g\exp Xh_0$，其$\mathfrak m$分量就是过渡函数；故光滑。在这些坐标中$\pi\Psi$是第一分量投影，所以是下沉，切空间为$\mathfrak g/\mathfrak h$。这些坐标也提供局部光滑截面；借截面检验可知任何使$\pi$为下沉的另一结构具有相同光滑映射，故结构唯一。左作用及正规商群的乘法/取逆均可借局部截面从$G$的光滑运算下降。正规性还给出$\operatorname{Ad}_G\mathfrak h=\mathfrak h$，微分后$\mathfrak h$为理想；$\dd\pi_e$保持括号并有核$\mathfrak h$，故得到商李代数。
\end{proof}

\begin{figure}[htbp]
\centering\DGDiagram{quotient-slice}

\caption{闭子群的局部乘积坐标。竖线的$X$固定、$Y$变化，落在同一左陪集；投影后仅留下横向的$\mathfrak m$坐标。图表示局部坐标而非群的全局直积。}
\label{fig:quotient-slice}
\end{figure}

\subsection{传递作用、各向同性与不变度量}
群作用约定为左作用$g\cdot(h\cdot p)=(gh)\cdot p$。固定$p$，稳定子$H=G_p$是闭子群，因为它是连续轨道映射$F_p(g)=g\cdot p$下$\{p\}$的原像。对$X\in\mathfrak g$定义基本场
\begin{equation}\label{eq:lie-fundamental-sign}
 X_M(q)=\left.\frac{\dd}{\dd t}\right|_0\exp(tX)\cdot q,
 \qquad [X_M,Y_M]=-[X,Y]_M.
\end{equation}
反号可由$F_p$证明：$X^R$与$X_M$相关，故\eqref{eq:lie-right-sign}传到轨道；对每个$p$如此即可。基本场的流是$\exp(tX)\cdot q$。若选择负指数定义基本场，则会得到同态约定；本书使用上面的正指数及反同态约定。

\begin{lemma}\label{lem:lie-surjective-rank}
第二可数流形之间的满射常秩光滑映射，其秩等于目标维数。
\end{lemma}
\begin{proof}
若秩$r$小于目标维数，常秩坐标把局部像放在一个$r$维坐标片内。用可数个紧集$K_j$覆盖源流形，每个$K_j$含于这样的源坐标邻域且其像含于对应目标坐标邻域；这是局部紧致性与第二可数性给出的可数预紧小邻域覆盖。每个$f(K_j)$紧、闭且无内点。取目标内一个闭坐标球，在其内部避开$f(K_1)$选更小闭球，再避开$f(K_2)$选下一球，依次令半径趋零，且下一球落在上一球内部。所有球都在第一个紧坐标球内，故交集非空；交中点不在任何$f(K_j)$，与满射矛盾。这给出了此处需要的可数纲论证，不调用Sard定理。
\end{proof}

\Needspace{18\baselineskip}
\begin{theorem}\label{thm:lie-homogeneous}
若$G$在$M$上光滑传递作用，则$gH\mapsto g\cdot p$给出$G/H\simeq M$的微分同胚，$H=G_p$。$G$不变黎曼度量与$\mathfrak g/\mathfrak h$上在各向同性表示
\[
 h\cdot(X+\mathfrak h)=\operatorname{Ad}_hX+\mathfrak h
\]
下不变的正定内积一一对应。特别地，$H$紧时存在这样的度量。
\end{theorem}
\begin{proof}
由$F_p\circ L_g=(g\cdot)\circ F_p$，轨道映射常秩；传递性及上引理使它为下沉。其单位元微分的核等于$\mathfrak h$：若$X_M(p)=0$，基本场的积分曲线从$p$出发只能是常值曲线，故$\exp(tX)\in H$；反向显然。于是诱导的双射$G/H\to M$微分处处为同构，逆函数定理使它成为局部微分同胚；双射再使局部逆拼成全局光滑逆。

不变度量在$p$处的内积必须被稳定子保持。对曲线$\exp(tX)H$作用$h\in H$，得$h\exp(tX)H=\exp(t\operatorname{Ad}_hX)H$，所以各向同性表示确如所写。反过来，把该内积用$g$的作用搬到$g\cdot p$；更换代表$g$为$gh$不改变结果，局部截面保证光滑性。紧$H$时用定理\ref{thm:lie-compact-average}对这个商空间表示平均。
\end{proof}

\begin{example}[球面的齐性表示]
$\mathrm{SO}(n+1)$在单位球面$S^n$上作用传递：将任意单位向量补成正向正交基，即构造把北极送到该向量的旋转。北极$e_{n+1}$的稳定子为$\operatorname{diag}(Q,1)$，其中$Q\in\mathrm{SO}(n)$。于是$S^n\simeq\mathrm{SO}(n+1)/\mathrm{SO}(n)$。商切空间可用
$X_v=\left(\begin{smallmatrix}0&v\\-v^{\mathsf T}&0\end{smallmatrix}\right)$表示，$\dd F_e(X_v)=(v,0)$，稳定子作用为$v\mapsto Qv$，标准欧氏内积下降为单位球面的圆度量。
\end{example}

\subsection{测地线反射与对称空间}
\begin{definition}
连通黎曼流形$M$称为全局黎曼对称空间，若每个$p\in M$都存在全局等距映射$s_p:M\to M$，满足$s_p(p)=p$及$\dd(s_p)_p=-I$。若各点附近有这样一个局部等距映射，则称局部对称。
\end{definition}
等距映射保持Levi--Civita联络和测地线，故在法坐标小球内必有
\begin{equation}\label{eq:lie-geodesic-reflection}
 s_p(\exp_pv)=\exp_p(-v).
\end{equation}
这里$\exp_p$为黎曼指数，区别于$\exp_G$。双不变度量在$e$处使两者一致，一般情形不可混同。

\begin{lemma}\label{lem:lie-isometry-rigidity}
连通黎曼流形的等距映射由一点的像及该点的微分唯一确定。因而对称映射$s_p$唯一，且$s_p^2=\mathrm{id}$。
\end{lemma}
\begin{proof}
只需证明等距映射$f$若固定$p$且$\dd f_p=I$，便为恒等。它在$p$的法坐标邻域中由保持测地线而恒等。令$A=\{q:f(q)=q,\ \dd f_q=I\}$，同一论证使$A$开；连续性使其闭（在局部坐标中比较微分）。$A$非空且流形连通，故$A=M$。应用于两个映射的复合及$s_p^2$即可。
\end{proof}

\begin{theorem}\label{thm:lie-symmetric-complete}
全局黎曼对称空间测地完备，并且其等距变换群作为抽象群作用传递。它满足$\nabla R=0$。
\end{theorem}
\begin{proof}
设单位速极大测地线$\gamma:(a,b)\to M$有有限右端点$b$。取$t_0<b$且足够接近$b$，再取$\delta>0$使$[t_0-\delta,t_0]\subset(a,b)$而$t_0+\delta>b$。反射曲线
$\widetilde\gamma(t)=s_{\gamma(t_0)}\gamma(2t_0-t)$在$t\in[t_0,t_0+\delta]$有定义，是测地线，在$t_0$处与$\gamma$位置、速度相同。ODE唯一性使它与原线在重叠区间相同，遂延拓过$b$，矛盾。左端同理，故测地完备。

连通流形的任意两点可用有限条短测地线段连起来：从一点可这样到达的集合开，其补集也开。在每条短段中点$q$的对称$s_q$交换两端。有限次复合便把起点送到终点，故等距群传递；这一步只用群的作用，不声称已经证明全体等距群是李群。

令$\mathcal R(U,V,W,Z)=\langle R(U,V)W,Z\rangle$。等距映射保持它及其协变导数$\nabla\mathcal R$，后者为五阶协变张量。在$p$处$s_p$的微分在五个槽中均为$-I$，故$\nabla\mathcal R=-(\nabla\mathcal R)$，只能为零。度量平行使这等价于$\nabla R=0$。
\end{proof}

\begin{figure}[htbp]
\centering\DGDiagram{symmetric-reflection}

\caption{法坐标中的测地反射把同一测地线上等时的点交换；速度先由微分变号，再由参数反向恢复，因而可接出延长的测地线。}
\label{fig:symmetric-reflection}
\end{figure}

\begin{theorem}[局部对称的曲率判别]\label{thm:lie-local-symmetric}
黎曼流形局部对称当且仅当$\nabla R=0$。
\end{theorem}
\begin{proof}
必要性与上面的五阶张量证明相同。充分性取$p$的小法坐标球，在$T_pM$中令$K_v(w)=R_p(w,v)v$。沿$\gamma_v(t)=\exp_p(tv)$的平行移动记$P_v(t)$。由$\nabla R=0$，在这个平行标架中Jacobi方程化为常系数方程
\[
 j''+K_vj=0,\qquad j(0)=0,\quad j'(0)=w.
\]
解的收敛矩阵幂级数给出
\[
 (\dd\exp_p)_vw=P_v(1)S(K_v)w,\qquad
 S(K)=\sum_{r=0}^{\infty}\frac{(-1)^rK^r}{(2r+1)!}.
\]
这里用第12章的指数微分与Jacobi场关系，参数区间是$[0,1]$。平行移动保持内积，故法坐标拉回度量为
\[
 (\exp_p^*g)_v(w,z)=\langle S(K_v)w,S(K_v)z\rangle_p.
\]
$K_{-v}=K_v$，所以此式在$(v,w,z)\mapsto(-v,-w,-z)$下不变。由此$\exp_p\circ(-I)\circ\exp_p^{-1}$为局部等距映射，固定$p$且微分为$-I$，即局部对称。该论证不需要完备性，也没有将局部反射直接宣称为全局映射。
\end{proof}

\begin{example}[三类全局模型]
半径$a$的圆球面上，$s_p(q)=2\langle p,q\rangle p/a^2-q$是环境空间的正交变换，固定$p$并在$T_pS^n$上等于$-I$，所以球面为对称空间。它不一定属于所选的$\mathrm{SO}(n+1)$传递群，这不影响对称性的定义。

平坦环面$\R^n/\Lambda$上，$s_{[p]}([q])=[2p-q]$与代表选择无关，是全局等距反射。一般平坦流形是否有全局反射还须检查其商关系，不能仅由$R=0$断言全局对称。

具有双不变度量的连通李群上，取逆是等距映射：在左平凡化中，$\theta_{g^{-1}}\dd\iota_gv=-\operatorname{Ad}_g\theta_gv$，而$\operatorname{Ad}_g$正交。因此$s_p(g)=pg^{-1}p$是全局等距映射，固定$p$且微分为$-I$。这也由定理\ref{thm:lie-symmetric-complete}重新得到其测地完备性。
\end{example}

至此我们证明了闭子群商、传递作用与不变度量、双不变群的曲率以及局部对称判别。把所有全局对称空间表示为特定李群的商，仍须全体等距群的Lie群结构及进一步结构理论；本章没有用这些尚未证明的定理来封闭论证。


\section{李群与齐性空间习题及详细解答}
15.3改述Kirillov讲义原题6.2的实数$n=2$情形；15.4合并原题2.2(1)、2.3并加入圆周计算；15.7改述原题2.5并补图坐标。该讲义原题的印刷/PDF页分别为82、18、19，详见来源表。其余五题为原创补充。本轮未核实梁书中对应的李群原题，不编造题号；基本群相关的原题2.2(2)未纳入，避免预用尚未建立的普遍覆盖群理论。

\Needspace{12\baselineskip}
\subsection*{习题15.1：仿射群、左右场与结构方程}
令$G=\{(a,b):a>0\}$，表示$x\mapsto ax+b$。求群运算、左右不变场、括号、指数及左Maurer--Cartan形式，并说明它不具有双不变黎曼度量。
\begin{proof}[解答]
复合给出$(a,b)(c,d)=(ac,b+ad)$，单位元$(1,0)$，逆为$(a^{-1},-b/a)$。以单位元处的$a,b$方向为$H,E$，微分左右乘法得
\[
 H^L=a\partial_a,\quad E^L=a\partial_b,\qquad
 H^R=a\partial_a+b\partial_b,\quad E^R=\partial_b.
\]
因此$[H^L,E^L]=E^L$而$[H^R,E^R]=-E^R$，即$[H,E]=E$。从$(1,0)$出发，$xH+yE$的积分曲线满足$a'=xa$、$b'=ya$，所以
\[
 \exp(xH+yE)=\left(e^x,\ y\frac{e^x-1}{x}\right),
\]
在$x=0$时商按连续值1解释。对偶于左不变基的一形式为$\theta^H=\dd a/a$、$\theta^E=\dd b/a$。直接求导得$\dd\theta^H=0$及$\dd\theta^E=-\theta^H\wedge\theta^E$，与$[H,E]=E$的Maurer--Cartan式吻合。

若有双不变内积，\eqref{eq:lie-ad-skew}取$X=H,Y=Z=E$，得到$2\langle E,E\rangle=0$，违反正定性。相应地$\operatorname{ad}_H=\left(\begin{smallmatrix}0&0\\0&1\end{smallmatrix}\right)$、$\operatorname{ad}_E=\left(\begin{smallmatrix}0&0\\-1&0\end{smallmatrix}\right)$，故Killing矩阵为$\operatorname{diag}(1,0)$。群连通且可微分同胚于$\R^2$，仍不满足双不变正定度量的条件。
\end{proof}

\Needspace{12\baselineskip}
\subsection*{习题15.2：三维旋转的指数与曲率}
令$\widehat v(w)=v\times w$。推导Rodrigues公式，计算$\mathfrak{so}(3)$的Killing形式，并对内积$\langle\widehat v,\widehat w\rangle=v\cdot w$计算双不变度量的截面及Ricci曲率。
\begin{proof}[解答]
向量三重积使$\widehat v^2=vv^{\mathsf T}-r^2I$，其中$r=\|v\|$，故$\widehat v^3=-r^2\widehat v$。把指数级数按奇偶及这个关系归并，得到
\[
 \exp(\widehat v)=I+\frac{\sin r}{r}\widehat v+
 \frac{1-\cos r}{r^2}\widehat v^2.
\]
$r=0$时按级数连续延拓。它固定$v$轴，在$v^\perp$上为转角$r$的旋转。$\operatorname{Ad}_Q\widehat v=\widehat{Qv}$对$Q\in\mathrm{SO}(3)$成立，因为保向正交变换保持叉积，所以所给内积确实$\operatorname{Ad}$不变。

由$\widehat v\widehat w=wv^{\mathsf T}-(v\cdot w)I$，迹为$-2v\cdot w$，这也是$\operatorname{ad}$表示中的迹，故$B=-2\langle\ ,\ \rangle$。对正交单位$v,w$，$\|v\times w\|=1$，从定理\ref{thm:lie-biinvariant-curvature}得到每个截面$K=1/4$，$\operatorname{Ric}=\frac12g$，标量曲率为$3/2$。若把度量乘常数$c>0$，截面曲率变成$1/(4c)$；固定Lie括号并不意味着所有归一化下曲率都为1。
\end{proof}

\Needspace{12\baselineskip}
\subsection*{习题15.3：实$\mathfrak{sl}(2)$的Killing形式}
由Kirillov原题6.2取实数$n=2$情形。直接计算$H,E,F$基中的三个$\operatorname{ad}$矩阵及Killing矩阵，证明$B(X,Y)=4\operatorname{tr}(XY)$并判定符号。
\begin{proof}[解答]
矩阵的列表示被作用的基向量。由$[H,E]=2E,[H,F]=-2F,[E,F]=H$，得到
\[
 \operatorname{ad}_H=\begin{pmatrix}0&0&0\\0&2&0\\0&0&-2\end{pmatrix},\quad
 \operatorname{ad}_E=\begin{pmatrix}0&0&1\\-2&0&0\\0&0&0\end{pmatrix},\quad
 \operatorname{ad}_F=\begin{pmatrix}0&-1&0\\0&0&0\\2&0&0\end{pmatrix}.
\]
相乘后取迹给出
\[
 (B_{ij})=\begin{pmatrix}8&0&0\\0&0&4\\0&4&0\end{pmatrix}.
\]
$4\operatorname{tr}(X_iX_j)$在这组基上有同一矩阵，双线性推出所有$X,Y$的等式。换基$H,E+F,E-F$后为$\operatorname{diag}(8,8,-8)$，故符号为$(2,1)$，非退化却不定。若这个李代数有正定$\operatorname{ad}$不变内积，则$\operatorname{ad}_H$为斜伴随，但$\operatorname{ad}_HE=2E$蕴含$4\|E\|^2=0$，矛盾。因此$\mathrm{SL}(2,\R)$也没有双不变黎曼度量；不定的双线性型不属于本章正定黎曼结论的适用范围。
\end{proof}

\Needspace{12\baselineskip}
\subsection*{习题15.4：局部同构、离散核与圆周周期}
合并Kirillov原题2.2(1)、2.3，并扩展圆周例子。证明连通群的离散正规子群在中心内；证明微分同构的连通李群同态为覆盖。再求所有光滑群同态$F:S^1\to S^1$，说明李代数同态何时能下降到圆周。
\begin{proof}[解答]
对离散正规$N$及$n\in N$，$g\mapsto gng^{-1}$是从连通空间到离散空间的连续映射，恒为$n$。这独立于是否存在商流形。微分同构使同态在单位元附近为局部微分同胚，像是开子群，目标连通使它满射。取小$U$使$UU^{-1}\cap\ker F=\{e\}$，则定理\ref{prop:lie-covering}的证明给出$F^{-1}(F(U))=\coprod_{k\in\ker F}kU$，每片均微分同胚到$F(U)$，这正是覆盖条件。核离散、正规，故在中心内；只说“局部可逆”尚未写出覆盖所需的不交原像片。

圆周单位元切空间以$i$为基，设$\dd F_e(i)=ci$，$c\in\R$。指数自然性给出$F(e^{it})=e^{ict}$。因$e^{i(t+2\pi)}=e^{it}$，必须$e^{2\pi ic}=1$，所以$c=m\in\mathbb Z$。反过来$F(z)=z^m$定义光滑同态。$m=0$时为常同态；$m\ne0$时导数可逆，核有$|m|$个点，成为$|m|$重覆盖。因此所有实数$c$都给出一维交换李代数同态，但只有整数保持周期格而下降为圆周同态。
\end{proof}

\Needspace{12\baselineskip}
\subsection*{习题15.5：环面中的无理绕行}
取无理数$\alpha$，令$F:\R\to T^2$为$F(t)=(e^{it},e^{i\alpha t})$。证明它是单射浸入同态，像稠密但不为全体，且不是到其像（取子空间拓扑）的同胚。说明$T^2/F(\R)$为何不能用闭子群商定理。
\begin{proof}[解答]
同态性直接来自指数加法。若$F(t)=(1,1)$，则$t=2\pi m$且$\alpha m\in\mathbb Z$，无理性迫使$m=0$。导数为非零向量$(1,\alpha)$的平移，所以处处单射，故为浸入。

先证明$\{e^{2\pi i m\alpha}:m\in\mathbb Z\}$在$S^1$稠密。把$N+1$个小数部分$\{j\alpha\}$分入$N$个等长区间，必有两个相差小于$1/N$。取差并允许改变符号，得到整数$q$及$k$使$0<q\alpha-k=\delta<1/N$（此处$q$允许负值）。角度为$\delta$的非负整数倍在$[0,1]$中形成最大间隙不超过$\delta$的网格，每个网格点都属于该循环子群；令$N\to\infty$即可任意逼近。

对给定$(e^{i\theta},e^{i\phi})$，取$t=\theta+2\pi m$可固定第一坐标，第二坐标由上述稠密性逼近目标，故像在$T^2$稠密。但第一坐标为1时，像中的第二坐标只有可数多个；圆周不可数，所以像不为全体。

上述抽屉构造还给出非零整数$q_j$，使$e^{2\pi iq_j\alpha}\to1$。可选$|q_j|\to\infty$，否则某非零$q$重复出现便会使$q\alpha$为整数。于是$F(2\pi q_j)\to(1,1)$，而原像不趋于0，故逆不连续。最后商中单位陪集的原像$F(\R)$不闭，因此该点不闭，商甚至不是$T_1$空间。群的像为浸入子群与其为闭嵌入子群是不同结论。
\end{proof}

\Needspace{12\baselineskip}
\subsection*{习题15.6：单位分支内的非指数元素}
在$\mathrm{SL}(2,\R)$中令$A=\operatorname{diag}(-2,-1/2)$。证明它位于单位元的路径分支，却不等于任何实矩阵的指数。
\begin{proof}[解答]
记$R(\theta)=\left(\begin{smallmatrix}\cos\theta&-\sin\theta\\\sin\theta&\cos\theta\end{smallmatrix}\right)$，则
\[
 c(t)=R(\pi t)\operatorname{diag}(2^t,2^{-t}),\qquad0\le t\le1
\]
是行列式恒为1的光滑路径，连接$I$与$A$。若$e^X=A$且$X$实，则$X$与其指数交换，因而$XA=AX$。$A$的两个对角元互异，比较交换等式的非对角元得$X_{12}=X_{21}=0$，所以$X$是实对角矩阵，其指数对角元均为正，矛盾。特别地它也不是任何$\mathfrak{sl}(2,\R)$元素的指数。这个例子说明连通单位分支由指数邻域的有限乘积生成，不意味着每个元素都是单个指数。
\end{proof}

\Needspace{12\baselineskip}
\subsection*{习题15.7：Grassmann流形的稳定子与图坐标}
改述Kirillov原题2.5并补充坐标验证。对$0<k<n$，把$k$维线性子空间全体表示为正交群的齐性空间，求维数及一个显式局部坐标，并构造不变黎曼度量。
\begin{proof}[解答]
令$P_0=\R^k\oplus0$。任一$k$平面取正交归一基，再补为$\R^n$的正交基，故$\mathrm O(n)$作用传递。保持$P_0$的正交矩阵也保持$P_0^\perp$，恰为块对角$\mathrm O(k)\times\mathrm O(n-k)$。这是闭子群，所以可定义
\[
 \operatorname{Gr}(k,n)=\mathrm O(n)/(\mathrm O(k)\times\mathrm O(n-k)),
 \quad \dim\operatorname{Gr}(k,n)=k(n-k).
\]
横截于$P_0^\perp$的平面唯一写成$\{(x,Zx):x\in\R^k\}$，其中$Z\in M_{n-k,k}(\R)$。在$Z=0$附近，对列$(e_i,Ze_i)$及剩余标准基作Gram--Schmidt，分母均非零，所得正交矩阵$Q(Z)$光滑，且$Q(Z)P_0$恰为该图平面。反向若$Q$的左上$k\times k$块$A$可逆、左下块为$C$，则图坐标为$Z=CA^{-1}$。两个构造均光滑，故图坐标与商流形结构相容；在任意参考平面作同样构造得到全体坐标。

因此$T_{P_0}\operatorname{Gr}(k,n)\simeq\operatorname{Hom}(P_0,P_0^\perp)$。稳定子$(U,V)$在该空间作用为$Z\mapsto VZU^{-1}$，保持$\operatorname{tr}(Z^{\mathsf T}W)$。由定理\ref{thm:lie-homogeneous}它唯一延拓为正交群不变黎曼度量。这里讨论的是无定向平面；若换成定向Grassmann流形，稳定子须相应改变，不能原封不动沿用此商。
\end{proof}

\Needspace{12\baselineskip}
\subsection*{习题15.8：左不变度量下的两种指数}
取$\mathfrak{so}(3)$基$A,B,C$，满足$[A,B]=C,[B,C]=A,[C,A]=B$。规定它们正交且平方长度为$I_1,I_2,I_3>0$，左平移延拓度量。求$\exp_G(t(A+B))$的协变加速度，判断何时它是测地线，并解释此处与双不变公式的区别。
\begin{proof}[解答]
对任意左不变场$X,Z$，Koszul公式在常内积下给出
\[
 \langle\nabla_XX,Z\rangle=-\langle[X,Z],X\rangle.
\]
令$X=A+B$。$[X,A]=-C,[X,B]=C,[X,C]=A-B$，所以前两个配对为零，第三个为$-(I_1-I_2)$。因此
\[
 \nabla_XX=\frac{I_2-I_1}{I_3}C.
\]
群指数曲线的速度是$X^L$，沿线加速度就是该左不变场的限制。于是此曲线为测地线当且仅当$I_1=I_2$。当$I_1\ne I_2$时，$\exp_G(tX)$与相同初速度的黎曼测地线不同；即使$I_1=I_2\ne I_3$使这一个方向恰好为测地线，也不能推出度量双不变。双不变要求全部$\operatorname{ad}$斜伴随，依次代入三个基即迫使$I_1=I_2=I_3$。所以不能把群指数的一参数子群定义自动解释为任意左不变度量的黎曼指数。
\end{proof}

